\section*{Epistemolojik Manifest: Kuramsal Paradigma Değişimi ve Doğa Felsefesi}
\addcontentsline{toc}{section}{Epistemolojik Manifest: Kuramsal Paradigma Değişimi ve Doğa Felsefesi}
\label{sec:epistemolojik_manifest}

\subsection*{1. Epistemolojik Sadelik ve Konformal Geometrik Perspektif}

Modern teorik fizikte ve gözlemsel kozmolojide karşılaşılan temel açık sorular—karanlık madde, karanlık enerji, kozmik şişme ve kuantum ölçekleri ile kütleçekim arasındaki kavramsal uyum—farklı kuramsal ve geometrik çerçevelerin araştırılmasını gerektirmektedir:
\begin{itemize}
    \item Galaksi dönüş eğrilerinin Keplerian dinamikten sapması, soğuk karanlık madde hipotezinin geliştirilmesine yol açmıştır.
    \item Tip Ia süpernovalarının parlaklık-kırmızıya kayma bağıntısı, kozmik ivmelenme ve karanlık enerji varsayımını doğurmuştur.
    \item Ufuk ve düzlük problemleri, erken evren kozmik enflasyon modellerini gündeme getirmiştir.
    \item Hadronların hapsi ve kütle kazanımı mekanizmaları ise perturbative olmayan alan teorilerinde araştırılmaya devam etmektedir.
\end{itemize}

Bu araştırma monografisi, metrik uzay genişlemesi yerine durağan bir arka planda maddesel bağlı yapıların yerel konformal ölçek büzülmesini ($r_{\mathrm{atom}}(t) \to 0$) ele alarak, bu gözlemsel olguların konformal simetri ve büzülen ölçüm standartları altında nasıl tezahür ettiğini incelemektedir. Bu yaklaşım, bilim felsefesinde Occam'ın usturası ilkesi doğrultusunda, asgari aksiyomatik varsayımla geniş bir gözlemsel spektrumun açıklanabilirliğini test etmeyi amaçlayan kuramsal bir araştırma programıdır.

\subsection*{2. Matematiksel Meşruiyet: Konformal İkizlik ve Nilpotent Cebir}

Uzayın genişlemesi ile atomik cetvellerin büzülmesi arasındaki ilişki bir soyutlama tercihinden ibaret kalmayıp, diferansiyel geometrinin kesin bir teoremine dayanır. Genel görelilikte metrik tensörün konformal dönüşümü:
\begin{equation}
    g_{\mu\nu} \to \Omega^2(t) g_{\mu\nu}
\end{equation}
koordinat sistemleri arasında tam bir geometrik eşdeğerlik (konformal ikizlik / conformal duality) kurar. Christof Wetterich (2013), Fred Hoyle, Jayant Narlikar ve Roger Penrose gibi kuramcıların farklı bağlamlarda işaret ettiği üzere; metrik uzayın genişlemesi ($a(t) \propto e^{H_0 t}$) ile maddesel uzunluk cetvellerinin daralması ($\lambda(t) \propto e^{-H_0 t}$) matematiksel olarak özdeş fiziksel öngörüler üretir.

Bu çalışmada geliştirilen $\kappa^2 = 0$ dual nilpotent cebir ve $\mathbb{D}_t \equiv \frac{d}{dt} + \kappa H_0$ Ölçek-Kovaryant Türevi, söz konusu konformal simetriyi pasif bir koordinat seçimi seviyesinden çıkararak nedensel ve dinamik bir fiziksel ontolojiye dönüştürmüştür. Uzay durağan Öklid sürekliliğini korurken, madde kendi içsel fazı boyunca büzülen dinamik bir Skalor olarak tanımlanmıştır.

\subsection*{3. Temel Sabitlerin Fiziksel Anlamı: $c$ ve Sommerfeld Sabiti $\alpha$}

Standart fizik, doğadaki temel boyutsuz ve boyutlu sabitlerin kökeni karşısında açıklamasız kalmış; bu değerleri evrenin başlangıcında belirlenmiş keyfi doğa sabitleri olarak kabul etmiştir:
\begin{itemize}
    \item \textit{``Işık hızı neden yaklaşık $3 \times 10^8\text{ m/s}$'dir ve neden doğanın aşılamaz mutlak tavanıdır?''}
    \item \textit{``Sommerfeld ince yapı sabiti $\alpha$'nın atomik kinematikteki fiziksel karşılığı nedir?''}
\end{itemize}

Bu çalışmada geliştirilen hidro-geometrik kalkülüs çerçevesinde:
\begin{enumerate}
    \item \textbf{Işık Hızı ($c$):} Harici bir hız limiti varsayımı olmaktan çıkarak, durağan elastik vakum dokusunun enine faz iletim hızı ($v_{\mathrm{faz}} = 1/\sqrt{\varepsilon_0\mu_0} = c$) ve bir vakum girdabı olan maddenin karakteristik radyal büzülme tavan hızı ($v_c = c$) olarak analitik temele kavuşturulmuştur.
    \item \textbf{İnce Yapı Sabiti ($\alpha \approx 1/137.036$):} Çekirdeğin içsel radyal büzülme hızı ($v_r = v_c = c$) karşısında, elektronun ışımasal çöküşe uğramadan kararlı bir yörüngede dolanabilmesi için çizdiği \textbf{geometrik spiral hatve oranıdır (vortex pitch ratio)}:
          \begin{equation}
              \alpha \equiv \frac{v_\theta}{v_r} = \frac{v_{\mathrm{Bohr}}}{c} = \frac{e^2}{4\pi\varepsilon_0\hbar c} \approx \frac{1}{137.036}.
          \end{equation}
          $Z\alpha \to 1$ eşiği ise teğetsel hızın büzülme tavanına ($v_\theta = c$) dayandığı ve süperkritik Dirac alanının vakum kutuplanmasıyla ($e^+ e^-$ çift üretimi) sistemi kararlı sınıra kilitlediği alt-atomik olay ufkudur.
\end{enumerate}
Bu iki temel sabitin hidrodinamik ve topolojik birer orana indirgenmesi, teorik fiziğin bir asırlık ontolojik kördüğümünü çözmektedir.

\subsection*{4. Boyutsuz Oranların Değişmezliği Kalkanı: Yerel Laboratuvar Testleri (Eötvös / MICROSCOPE)}

Birleşik Büzülme Hipotezi'nin yerel laboratuvar deneyleriyle uyumu, boyutsuz oranların korunumu ilkesiyle teminat altına alınmıştır. Eğer atomik cetveller kütle ölçeklemesinden yalıtılmış olarak tekil biçimde büzülseydi, MICROSCOPE uydu deneyi ve Eötvös burulma terazileri diferansiyel ivme sapmaları kaydederek teoriyi doğrudan yanlışlardı.

Bağlı tüm maddesel durumlar evrensel ölçek kuralı gereğince eş zamanlı olarak $m(t) = m_0 e^{H_0 t}$ konformal kütle kazandığından, laboratuvarda ölçülen herhangi iki kütlenin boyutsuz oranı mutlak olarak korunur:
\begin{equation}
    \frac{d}{dt} \left( \frac{m_{\mathrm{test}}}{m_{\mathrm{ref}}} \right) = \frac{d}{dt} \left( \frac{m_{1,0} e^{H_0 t}}{m_{2,0} e^{H_0 t}} \right) \equiv 0.
\end{equation}
Benzer biçimde İnce Yapı Sabiti $\alpha$ ve proton-elektron kütle oranı $\mu = m_p/m_e$ zamandan bağımsız kalır ($\dot{\alpha}/\alpha = 0, \dot{\mu}/\mu = 0$). Bu korunum, yerel Zayıf Eşdeğerlik İlkesi (WEP) testlerinin $\eta < 10^{-15}$ sınırında neden sıfır diferansiyel ivme ölçtüğünü tam bir tutarlılıkla açıklar.

\subsection*{5. Gökyüzü ile Yerkürenin Birleşimi: Yerel Ayrışma, Efemeris Sınırları ve Paleojeoloji}

Kozmolojik modeller sıklıkla milyarlarca ışık yılı uzaktaki soyut parametrelerle sınırlı kalarak yerel denetimden uzaklaşır. Birleşik Büzülme Hipotezi, derin uzay astrofiziği ile Güneş Sistemi yerel dinamiklerini tutarlı bir çerçevede birleştirir:
\begin{itemize}
    \item \textbf{Yerel Virial Ayrışma ve LLR Ölçümleri:} Kütleçekimsel potansiyelin baskın olduğu bağlı sistemlerde ($\nabla^2\Phi \gg 3H_0^2$), yerel yörüngeler McVittie kuralı uyarınca kozmik büzülme akışından virial olarak ayrışır. Apollo retroreflektörleriyle ölçülen Ay'ın $3.82\text{ cm/yıl}$'lık açılması, okyanus altimetrisiyle (TOPEX/Poseidon) doğrulanan gelgit sürtünmesi mekanizmasıyla tam bir uyum içindedir.
    \item \textbf{Efemeris ve Radar Sınırları:} Atomik saniyelerin büzülme ritmi yerel koordinatları değil atomik frekansları tanımladığından, gezegen efemerisleri ve radar telemetrisi (1 AU için $< 1.5\text{ cm/yıl}$) ile mutlak bir deneysel uyum sağlanır; kuram yerel gök mekaniği testlerinden başarıyla geçer.
\end{itemize}

\subsection*{6. Kavramsal Netlik: ``Vakum Akışkanı'' ve 19. Yüzyıl Eter Algısının Aşılması (Türeyen Uzay-Zaman)}

Modelde kullanılan ``vakum akışı'', ``girdap'' ve ``Bernoulli basıncı'' kavramları, 19. yüzyılın klasik mekanik eter hipotezleriyle karıştırılmamalıdır. Klasik eter teorileri uzayı sürtünmeli, mekanik dişlilerden oluşan maddesel bir sıvı olarak tasavvur etmiş ve Michelson--Morley deneyleriyle terk edilmiştir.

Bu çalışmada arka plan sürekliliği $\mathcal{E}_\infty$, sürtünmesiz, eylemsiz, kütlesiz ve mutlak pürüzsüz bir Öklid uzayıdır. Maddenin büzülmesiyle oluşan vakum akışı, modern kuramsal fizikteki \textit{Türeyen Uzay-Zaman (Emergent Spacetime)} ve \textit{Akustik/Hidrodinamik Yerçekimi (Analogue Gravity)} yaklaşımlarıyla tam bir metodolojik uyum sergiler. Parçacıklar durağan bir zemine gömülü harici bilyeler olarak değil; arka plan geriliminin topolojik girdap düğümleri olarak varlık kazanır.

\subsection*{7. Hesaplamalı Doğrulama ve Açık Kaynak Çözücüsü: Minimal CMB Kütüphanesi ve Sayısal Yol Haritası}

Birleşik Büzülme Hipotezi kavramsal ve analitik çerçevesini tamamlamakla kalmamış; Planck 2018 CMB tam açısal güç spektrumunu ($\ell = 2\text{--}2500$) fit eden açık kaynaklı, bağımsız araştırmacılarca tekrarlanabilir ve DOI tescilli minimal bir Python hesaplama çözücüsü (\texttt{scale-cmb-solver}, \href{https://github.com/FurkanHaydari/universe-is-not-expanding}{\texttt{FurkanHaydari/universe-is-not-expanding}}) ile somutlaştırılmıştır:
\begin{enumerate}
    \item \textbf{Planck CMB Açısal Güç Spektrumu Sayısal Çözücüsü:} Çalışmamızın minimal Python çözücüsü (\texttt{analysis/cmb\_spectrum\_analysis.py}); büzülen plazmanın duran dalga transfer denklemi, baryon yüklemesi ($R_b \approx 0.058$), akustik ölçek periyodu ($\Delta\ell \approx 293.5$), erken ISW potansiyel desteği ve Silk difüzyon sönümlemesi ($\ell_D \approx 1638.7$) mekanizmalarıyla Planck 2018 sıcaklık anizotropisi ($D_\ell^{\mathrm{TT}}$) verilerini (39 binned veri noktası) doğrudan modellemektedir. Çözücü; 1. Akustik Tepeyi $\ell_1 = 220.9$ (gözlem: 220.0), 2. Tepeyi $\ell_2 = 518.4$ (gözlem: 537.5), 3. Tepeyi $\ell_3 = 820.2$ (gözlem: 810.8) ve tepe oranlarını ($H_1/H_2 = 2.169$, $H_3/H_2 = 0.886$) Planck verileriyle tam bir ampirik uyumla üretmektedir. Bu minimal açık kaynak çekirdeği, standart CAMB ve CLASS algoritmalarının büzülen cetvel formatındaki tam Boltzmann hiyerarşisi simülatörüne genişletilmektedir.
    \item \textbf{Döteryum ve İlksel Bollukların Nükleer Ağ Simülasyonu:} Galaksiler arası gazda gözlenen $D/H \sim 2.5 \times 10^{-5}$ bolluk dengesi; yıldız içi astrasyon (döteryumun yanarak tükenmesi) ile durağan vakum akısı kaynaklı fotonükleer üretim mekanizmasının dengesini kuran çok bileşenli bir \textit{Astrofiziksel Nükleer Ağ Simülatörü} ile sayısal doğrulamaya kavuşturulmaktadır.
\end{enumerate}
Bu açık kaynaklı yazılım paketi, analitik mimarinin ampirik Planck ve Pantheon+ verileriyle doğrudan örtüştüğünü bilgisayar kodu seviyesinde tescil etmektedir.

\subsection*{8. Bilim Sosyolojisi ve Kurumsal Kabul Direnci}

Tarih boyunca bilimsel paradigmaların dönüşümü, Max Planck'ın işaret ettiği sosyolojik gerçeklikle ilerlemiştir: \textit{``Yeni bir bilimsel gerçek, muhaliflerin zamanla sahneden çekilmesi ve yeni neslin bu gerçeğe aşina olarak büyümesiyle zafer kazanır.''}

Bu çalışmanın muhatap olacağı asıl mukavemet, doğrudan kurumsal müesses nizamın (establishment) epistemolojik alışkanlıklarından kaynaklanacaktır:
\begin{itemize}
    \item Karanlık madde ve karanlık enerji varsayımlarının ortadan kalkması, bu alanlara ayrılmış devasa fonları ve onlarca yıllık kurumsal araştırma programlarını tartışmaya açacaktır.
    \item Tekil bir Büyük Patlama ($t=0$) tekilliğinin yerini maddenin büzülme fazı eşiğine bırakması, genişleyen metrik üzerine inşa edilmiş literatürün yeniden değerlendirilmesini gerektirecektir.
\end{itemize}
Bu nedenle çalışmamızda savunmacı bir polemik dili bütünüyle terk edilmiş; doğrudan tanımlar, analitik bağıntılar ve yanlışlanabilir somut öngörülerden oluşan aksiyomatik bir fizik mimarisi inşa edilmiştir. Doğa felsefesinin en yalın haliyle ortaya konan bu kuram, matematiksel ve gözlemsel testler karşısında tam bir tutarlılığa sahiptir.

\vspace{1.5em}
\hrule
\vspace{1.5em}
